\section{Experiments}
\label{sec:experiments}

\subsection{Setup}

\paragraph{Datasets.}
We use three HSI+LiDAR benchmarks: Trento (63 HSI bands, 1 LiDAR channel, 6 classes), Houston 2013 (144 bands, 1 LiDAR, 15 classes), and MUUFL Gulfport (72 bands, 2 LiDAR channels, 11 classes). HSI is resampled to 36 unified bands; LiDAR is padded to 2 channels. Patch size is $7\times7$.

\paragraph{Marathon CIL protocol.}
Default ordering: Trento (tasks 0--2) $\to$ Houston (tasks 3--5) $\to$ MUUFL (tasks 6--8). Each task adds 2--5 new classes. We evaluate task-agnostic accuracy (TAg): balanced accuracy over all seen classes without task identifiers.

\paragraph{Implementation.}
\backbonename{} backbone uses embed\_dim=64, 4 spatial VSSBlocks, 2 spectral Mamba layers. Pre-trained on Task~0 (Trento, 2 classes). Warm-up: 3 tasks, 30 epochs, lr=$10^{-3}$. LoRA: rank $r{=}4$ (HSI), $r{=}8$ (LiDAR), 50 epochs, lr=$5{\times}10^{-4}$. Batch size 128. No orthogonal constraint, no KD loss (ablation shows these do not help). Results averaged over 2 seeds.

\paragraph{Baselines.}
All baselines use the same \backbonename{} backbone for fair comparison. Exemplar-free (EF): Frozen NCM, EWC~\cite{kirkpatrick2017ewc}, LwF~\cite{li2017lwf}, Analytic CL (CrossACL-style)~\cite{xu2025crossacl}, PSCEN~\cite{pscen2025}, FOSTER~\cite{wang2022foster}. Pseudo-replay (PR): DCPRN~\cite{deng2024dcprn}. Exemplar-based (EB, 5 samples/class): iCaRL~\cite{rebuffi2017icarl}, LUCIR~\cite{hou2019lucir}.

\subsection{Main Results}

Table~\ref{tab:main} presents the full Marathon CIL results. Among exemplar-free methods, \methodname{}+SHINE achieves the highest accuracy at 75.7\%, surpassing DCPRN+SHINE (74.6\%) by +1.1pp and Frozen+SHINE (72.7\%) by +3.0pp. The improvement is consistent across all three domains: +4.5pp on Trento, +1.4pp on Houston, and +4.2pp on MUUFL over the best frozen-backbone baseline.

Full backbone fine-tuning methods (EWC, LwF) perform \emph{worse} than frozen baselines under cross-domain CIL. EWC+SHINE (65.9\%) degrades 6.8pp below Frozen+SHINE, confirming that unrestricted parameter updates cause catastrophic forgetting across domain boundaries. \methodname{}'s selective freeze/adapt strategy avoids this pitfall.

Exemplar-based methods (iCaRL 82.1\%, LUCIR 83.0\%) outperform all exemplar-free methods, as expected. However, they require storing real samples from past domains---a constraint that may be impractical in remote sensing due to storage, privacy, or licensing restrictions.

\begin{table}[t]
\centering
\caption{Marathon CIL results on S\textsuperscript{2}CM backbone (9 tasks, 32 classes). Best exemplar-free result in \textbf{bold}. EF=Exemplar-Free, PR=Pseudo-Replay, EB=Exemplar-Based.}
\label{tab:main}
\begin{tabular}{llcccc}
\toprule
Method & Type & Trento & Houston & MUUFL & Avg TAg \\
\midrule
FOSTER+S & EF & 87.1 & 68.5 & 70.7 & 72.7 \\
EWC+S & EF & 84.8 & 63.0 & 59.4 & 65.9 \\
LwF+S & EF & 89.0 & 60.5 & 68.1 & 68.5 \\
Frozen+S & EF & 87.1 & 68.5 & 70.7 & 72.7 \\
Analytic+S & EF & 87.1 & 68.5 & 70.7 & 72.7 \\
PSCEN+S & EF & 87.1 & 68.5 & 70.7 & 72.7 \\
DCPRN+S & PR & 87.7 & 71.9 & 71.1 & 74.6 \\
\textbf{AnchorLoRA+S} & \textbf{EF} & \textbf{91.6} & \textbf{69.9} & \textbf{74.9} & \textbf{75.7} \\
\midrule
iCaRL+S & EB$^{5}$ & 92.0 & 81.9 & 76.9 & 82.1 \\
LUCIR+S & EB$^{5}$ & 96.5 & 84.7 & 73.5 & 83.0 \\
\bottomrule
\end{tabular}
\vspace{-1mm}
{\footnotesize +S denotes +SHINE domain normalization. EB$^5$: 5 exemplars per class.\\
Frozen, Analytic, PSCEN, and FOSTER+SHINE all converge to 72.7\% because SHINE normalization dominates the classifier choice when the backbone is frozen; the NCM, ridge regression, and calibration variants produce identical predictions after whitening.}
\end{table}

\subsection{Cross-Backbone Comparison}

To disentangle the contributions of the \backbonename{} backbone and the \methodname{} CIL strategy, we evaluate on a Dual-Stream LSGA-ViT backbone (Table~\ref{tab:crossbb}).

\begin{table}[t]
\centering
\caption{Cross-backbone comparison. \backbonename{}'s spectral-spatial separation enables \methodname{}'s asymmetric adaptation.}
\label{tab:crossbb}
\begin{tabular}{llc}
\toprule
Backbone & CIL Method & Avg TAg \\
\midrule
Dual-LSGAVIT & Frozen NCM & 66.6 \\
Dual-LSGAVIT & AnchorLoRA+SHINE & 46.7 \\
\midrule
\backbonename{} & Frozen NCM & 59.2 \\
\backbonename{} & Frozen+SHINE & 72.7 \\
\backbonename{} & \textbf{AnchorLoRA+SHINE} & \textbf{75.7} \\
\bottomrule
\end{tabular}
\end{table}

\methodname{} on the ViT backbone (46.7\%) underperforms frozen ViT (66.6\%), a result that \emph{validates} our core thesis: AnchorLoRA is a backbone-aware co-design whose effectiveness stems from \backbonename{}'s explicit spectral-spatial separation. Without distinct modality branches, LoRA cannot distinguish which features to preserve and which to adapt. Rather than a limitation, this result confirms that modality-aware architecture is essential for asymmetric adaptation---the method and backbone form a principled co-designed system.

\subsection{Domain Ordering Robustness}

Table~\ref{tab:ordering} shows \methodname{} performance across 3 representative domain orderings (out of 6 possible permutations). The method consistently improves over SHINE across all tested orderings (+4.4 to +8.3pp), confirming that the spectral anchor is domain-agnostic.

\begin{table}[t]
\centering
\caption{Ordering robustness. \methodname{} improves over SHINE baseline across 3 representative domain orderings.}
\label{tab:ordering}
\begin{tabular}{lcc}
\toprule
Ordering & Avg TAg & $\Delta$ SHINE \\
\midrule
Trento$\to$Houston$\to$MUUFL & 74.4 & +8.3 \\
Houston$\to$MUUFL$\to$Trento & 73.6 & +5.8 \\
MUUFL$\to$Trento$\to$Houston & 78.7 & +4.4 \\
\bottomrule
\end{tabular}
\vspace{-1mm}
{\footnotesize All with warmup=3, rank=4/8, $\lambda_{\text{ortho}}$=0.}
\end{table}

\subsection{Ablation Study}
\label{sec:ablation}

Table~\ref{tab:ablation} presents comprehensive ablation results, revealing several counter-intuitive findings.

\begin{table}[t]
\centering
\caption{Ablation study on the Marathon CIL benchmark (seed 0).}
\label{tab:ablation}
\begin{tabular}{lcc}
\toprule
Configuration & Avg TAg & $\Delta$ \\
\midrule
\textbf{Best config} (no ortho, r4/8, w3) & \textbf{75.4} & --- \\
+ Orthogonal constraint ($\lambda$=0.1) & 74.4 & $-$1.0 \\
+ Spectral KD loss & 75.2 & $-$0.2 \\
Symmetric rank (r4/r4) & 71.5 & $-$3.9 \\
Larger rank (r8/r16) & 74.5 & $-$0.9 \\
Warmup = 2 tasks & 74.3 & $-$1.1 \\
Domain-selective LoRA & 70.8 & $-$4.6 \\
DualSpeed dual-head & 72.3 & $-$3.1 \\
SpecRoute (spectral routing) & 72.1 & $-$3.3 \\
\bottomrule
\end{tabular}
\end{table}

\paragraph{Orthogonal constraints hurt.}
Adding InfLoRA-style orthogonal regularization reduces accuracy by 1.0pp. With rank-4 adapters, the subspace is too small for orthogonal constraints to be useful; they instead limit later tasks' capacity.

\paragraph{Spectral KD is ineffective.}
Distilling spectral knowledge into the full-branch classifier provides negligible benefit ($-$0.2pp), because the spectral branch is already frozen and its information is directly available at inference.

\paragraph{Asymmetric rank is critical.}
Switching to symmetric rank ($r$=4 for both modalities) loses 3.9pp. The LiDAR branch requires 2$\times$ capacity because its features drift 2.5$\times$ more than HSI-spatial features (Table~\ref{tab:drift}).

\paragraph{Routing and dual-head designs fail.}
Both spectral-routed expert selection (SpecRoute, $-$3.3pp) and dual-head calibrated blending (DualSpeed, $-$3.1pp) underperform simple LoRA accumulation. Routing fails because spectral centroids lack sufficient discriminability for fine-grained expert selection. Accumulated LoRA provides beneficial cross-domain feature diversity that routing disrupts.

\paragraph{Domain-selective LoRA hurts.}
Applying only same-domain LoRA at inference ($-$4.6pp) confirms that cross-domain adapter accumulation provides useful feature diversity rather than interference.
